% Agentic Coding for Economists - Padova five-hour workshop
% Build from this directory:
% latexmk -pdf -interaction=nonstopmode workshop_slides.tex

\input{../shared/beamer_preamble}

\renewcommand{\WorkshopRepoUrl}{\url{https://github.com/meleantonio/AC4E_Padova}}
\renewcommand{\ParentCourseUrl}{\url{https://github.com/meleantonio/AC4E_Pavia}}
\title{Agentic Coding for Economists}
\subtitle{Padova Workshop - Five-Hour Research Harness}
\author{Antonio Mele}
\institute{Agentic Coding for Economists}
\date{University of Padova, 25 June 2026}

\newcommand{\demo}[1]{\textcolor{courseteal}{\textbf{Demo: #1}}}
\newcommand{\handoff}[1]{\textcolor{courseblue}{\textbf{Handoff: #1}}}

\begin{document}

% =============================================================================
\section{Start}
% =============================================================================

\begin{frame}{}
  \begin{center}
    {\LARGE \textbf{Agentic Coding for Economists}} \\
    \vspace{0.4em}
    {\large Padova five-hour workshop} \\
    \vspace{0.6em}
    Antonio Mele \\
    \vspace{0.4em}
    {\small University of Padova, 25 June 2026} \\
    \vspace{0.8em}
    {\footnotesize Participant repo: \WorkshopRepoUrl}
  \end{center}
\end{frame}

\begin{frame}{How To Use This Session}
  \begin{block}{Pattern}
    Short explanation, bounded demo, participant action, verification note.
  \end{block}
  \begin{itemize}
    \item Keep \artifact{START\_HERE.md}, \artifact{SCHEDULE.md}, and \artifact{GUIDE.md} open.
    \item Work from one tool lane: Codex, Claude Code, or Cursor.
    \item Treat agents as implementation assistants, not research owners.
    \item Every delegated task ends with evidence: diff, command output, PR, or review log.
  \end{itemize}
\end{frame}

\begin{frame}{Five-Hour Arc}
  \begin{center}
    \begin{tikzpicture}[
      node distance=0.35cm,
      phase/.style={rectangle, draw, rounded corners, minimum height=0.75cm, align=center, font=\scriptsize},
      arr/.style={-Stealth, thick}
    ]
      \node[phase, fill=courseblue!15, minimum width=1.75cm] (a) {Orient};
      \node[phase, fill=courseteal!15, minimum width=1.75cm, right=of a] (b) {Example};
      \node[phase, fill=courseblue!15, minimum width=1.75cm, right=of b] (c) {Harness};
      \node[phase, fill=courseteal!15, minimum width=1.75cm, right=of c] (d) {Orchestrate};
      \node[phase, fill=courseblue!15, minimum width=1.75cm, right=of d] (e) {Sprint};
      \node[phase, fill=courseteal!15, minimum width=1.75cm, right=of e] (f) {Verify};
      \draw[arr] (a) -- (b) -- (c) -- (d) -- (e) -- (f);
    \end{tikzpicture}
  \end{center}
  \vspace{0.4em}
  \begin{itemize}
    \item 00:00-00:45: orientation, safety, agent loop.
    \item 00:45-01:20: Card-Krueger scaffold.
    \item 01:30-02:15: skills, subagents, hooks, MCP.
    \item 02:15-03:00: issues, cloud agents, swarms, loops, goals.
    \item 03:00-05:00: sprint, verification, adoption.
  \end{itemize}
\end{frame}

\begin{frame}{Participant Repository Path}
  \begin{columns}[T,onlytextwidth]
    \begin{column}{0.48\textwidth}
      \textbf{Read first}
      \begin{itemize}
        \item \artifact{README.md}
        \item \artifact{START\_HERE.md}
        \item \artifact{SCHEDULE.md}
        \item \artifact{GUIDE.md}
        \item \artifact{docs/sources.md}
      \end{itemize}
    \end{column}
    \begin{column}{0.48\textwidth}
      \textbf{Use during demos}
      \begin{itemize}
        \item \artifact{examples/card-krueger/}
        \item \artifact{examples/starter\_article/}
        \item \artifact{materials/tool\_tracks.md}
        \item \artifact{templates/}
        \item \artifact{slides/workshop/demo\_map.md}
      \end{itemize}
    \end{column}
  \end{columns}
\end{frame}

\begin{frame}{What You Produce}
  \begin{itemize}
    \item A scoped project brief or design memo.
    \item Agent instructions and privacy boundaries.
    \item One usable harness piece: skill, subagent, hook, MCP config, goal, or loop.
    \item One GitHub issue or cloud-agent task with acceptance criteria.
    \item One verification artifact: test output, review log, table check, or screenshot.
    \item A seven-day adoption plan for one real research project.
  \end{itemize}
\end{frame}

\begin{frame}{Safety Boundaries}
  \begin{block}{Human-owned decisions}
    Research question, identification, data rights, confidentiality,
    interpretation, and final claims stay with the researcher.
  \end{block}
  \begin{itemize}
    \item Do not upload confidential referee material, private data, API keys, or student records.
    \item Use synthetic or public examples for workshop demos.
    \item Start read-only; allow edits only after scope is explicit.
    \item Require commands and outputs before accepting generated code or claims.
  \end{itemize}
\end{frame}

\begin{frame}[fragile]{First Read-Only Prompt}
  \demo{orientation without edits}
  \begin{lstlisting}[style=markdown]
Read README.md, START_HERE.md, SCHEDULE.md, GUIDE.md,
examples/card-krueger/README.md, and
examples/starter_article/README.md. Do not edit files.

Explain:
1. the workshop path;
2. the Card-Krueger running example;
3. which files I should copy or personalize;
4. which files or folders I should not expose to agents.
  \end{lstlisting}
\end{frame}

% =============================================================================
\section{Running Case}
% =============================================================================

\begin{frame}{Running Case: Card And Krueger}
  \begin{block}{Question}
    What happened to fast-food employment in New Jersey relative to eastern
    Pennsylvania after New Jersey raised its minimum wage in April 1992?
  \end{block}
  \begin{itemize}
    \item Treatment group: New Jersey fast-food stores.
    \item Comparison group: eastern Pennsylvania fast-food stores.
    \item Timing: before and after the April 1992 policy change.
    \item Outcome: full-time-equivalent employment.
  \end{itemize}
\end{frame}

\begin{frame}{Public Sources And Teaching Data}
  \begin{itemize}
    \item Public source links live in \artifact{docs/sources.md}.
    \item David Card's page lists the associated public data archive and readme.
    \item NBER Working Paper 4509 gives the core research context.
    \item \artifact{examples/card-krueger/} uses synthetic teaching data so everyone can run the demo immediately.
  \end{itemize}
  \vspace{0.4em}
  \begin{block}{Caveat}
    The bundled CSV is not Card and Krueger's raw data and does not support a
    substantive causal claim.
  \end{block}
\end{frame}

\begin{frame}{Data Map Checklist}
  \begin{columns}[T,onlytextwidth]
    \begin{column}{0.52\textwidth}
      \textbf{Required fields}
      \begin{itemize}
        \item source and access note;
        \item unit of observation;
        \item treatment and comparison groups;
        \item outcome and units;
        \item timing and sample restrictions;
        \item transformations and caveats.
      \end{itemize}
    \end{column}
    \begin{column}{0.44\textwidth}
      \textbf{Workshop file}
      \begin{itemize}
        \item \artifact{examples/card-krueger/}
        \item \artifact{docs/data\_source\_map.md}
        \item \artifact{templates/data\_source\_map.md}
      \end{itemize}
    \end{column}
  \end{columns}
\end{frame}

\begin{frame}{Baseline DID Mechanics}
  \begin{block}{Teaching estimand}
    \[
      (\overline{Y}_{NJ,after} - \overline{Y}_{NJ,before})
      -
      (\overline{Y}_{PA,after} - \overline{Y}_{PA,before})
    \]
  \end{block}
  \begin{itemize}
    \item The script computes group means and the simple before/after comparison.
    \item Tests check the input file, group labels, sample counts, command path, and caveats.
    \item This is a workflow demo, not a replication claim.
  \end{itemize}
\end{frame}

\begin{frame}[fragile]{Run The Example}
  \demo{baseline command}
  \begin{lstlisting}[style=bash]
python3 -m pip install -r examples/card-krueger/requirements.txt
python3 examples/card-krueger/src/did_analysis.py
python3 -m pytest examples/card-krueger/tests
  \end{lstlisting}
  \vspace{0.3em}
  Expected headline:
  \begin{lstlisting}[style=markdown]
Toy DID estimate: 3.0 FTE workers.
Synthetic teaching data: not Card and Krueger's raw data.
  \end{lstlisting}
\end{frame}

\begin{frame}{From Example To Your Project}
  \begin{itemize}
    \item Use \artifact{examples/card-krueger/} as the shared reference during demos.
    \item Use \artifact{examples/starter\_article/} as the personalizable harness skeleton.
    \item Replace the teaching data map with your own data source map.
    \item Keep the same discipline: scope, evidence, review.
  \end{itemize}
\end{frame}

% =============================================================================
\section{Harness}
% =============================================================================

\begin{frame}{Agent Loop}
  \begin{center}
    \begin{tikzpicture}[
      node distance=1.0cm,
      loopbox/.style={rectangle, draw, rounded corners, minimum width=2.2cm, minimum height=0.7cm, align=center, font=\small},
      arr/.style={-Stealth, thick}
    ]
      \node[loopbox, fill=courseblue!15] (read) {Read\\context};
      \node[loopbox, right=of read] (plan) {Plan};
      \node[loopbox, right=of plan] (act) {Act\\tools};
      \node[loopbox, right=of act] (verify) {Verify};
      \draw[arr] (read) -- (plan) -- (act) -- (verify);
      \draw[arr] (verify.north) .. controls +(0,0.55) and +(0,0.55) .. (read.north);
    \end{tikzpicture}
  \end{center}
  \begin{itemize}
    \item Better context improves plans.
    \item Better acceptance criteria improves edits.
    \item Better verification improves trust.
  \end{itemize}
\end{frame}

\begin{frame}{Harness Pieces}
  \begin{tabular}{p{0.25\linewidth}p{0.66\linewidth}}
    \toprule
    Piece & Economist use \\
    \midrule
    Instructions & persistent project map and boundaries \\
    Skills & repeatable procedures for replication, review, writing \\
    Subagents & specialist reviewer contexts \\
    Hooks & deterministic reminders after edits or before risky commands \\
    MCP & structured access to external tools and public data \\
    Goals/loops & longer tasks with stop conditions and evidence \\
    \bottomrule
  \end{tabular}
\end{frame}

\begin{frame}{Project Instructions}
  \begin{itemize}
    \item Root \artifact{AGENTS.md} tells agents how to work in this repo.
    \item Project copies can use \artifact{AGENTS.md}, \artifact{CLAUDE.md}, Cursor rules, or all three.
    \item Include repo map, allowed files, verification commands, and do-not rules.
    \item Keep instructions short enough that participants can audit them.
  \end{itemize}
  \handoff{\artifact{templates/AGENTS\_economics.md}}
\end{frame}

\begin{frame}{Brief, Memo, Spec}
  \begin{itemize}
    \item \artifact{docs/project\_brief.md}: what the project is and is not.
    \item \artifact{docs/research\_design\_memo.md}: data, design, risks, and claims.
    \item \artifact{spec/intent.md}: optional spec-driven development path for larger projects.
    \item Agents can draft; humans own identification and interpretation.
  \end{itemize}
  \handoff{\artifact{templates/project\_brief.md}, \artifact{templates/research\_design\_memo.md}}
\end{frame}

\begin{frame}{Skills}
  \begin{block}{Use when the task is repeatable}
    A skill packages a procedure that should be invoked consistently.
  \end{block}
  \begin{itemize}
    \item replication checker;
    \item referee checklist;
    \item paper polisher;
    \item research SDD plan.
  \end{itemize}
  \handoff{\artifact{examples/starter\_article/.cursor/skills/}}
\end{frame}

\begin{frame}{Subagents And Reviewers}
  \begin{itemize}
    \item Use a separate reviewer role for identification, data, literature, or verification.
    \item Give the reviewer a narrow file set and a blockers-first output format.
    \item Reviewer agents should not silently fix their own findings.
    \item In GitHub, request a review after the implementation PR exists.
  \end{itemize}
  \handoff{\artifact{examples/starter\_article/.cursor/agents/}}
\end{frame}

\begin{frame}{Hooks}
  \begin{itemize}
    \item Hooks are deterministic reminders or commands around agent tool use.
    \item Example: after editing the DiD script, remind the agent to run pytest.
    \item Keep hook examples safe and explicit; participants trust them before use.
    \item If a tool version changes hook behavior, check official docs first.
  \end{itemize}
  \handoff{\artifact{examples/hooks/README.md}}
\end{frame}

\begin{frame}{MCP And Public Data}
  \begin{itemize}
    \item MCP gives an agent structured access to tools such as public-data APIs.
    \item Use read-only public sources first.
    \item Put secrets in environment variables, never in the repo.
    \item For Padova, FRED is the shared optional public-data example.
  \end{itemize}
  \handoff{\artifact{examples/cursor/.cursor/mcp.json.example}}
\end{frame}

\begin{frame}{Goals, Loops, And Stop Conditions}
  \begin{itemize}
    \item Use a goal for work that spans several steps.
    \item Define done criteria before starting.
    \item Define stop conditions: missing data, failed tests, ambiguous identification, tool limits.
    \item Record each loop in a review or orchestration log.
  \end{itemize}
  \handoff{\artifact{templates/long\_running\_goal\_prompt.md}}
\end{frame}

% =============================================================================
\section{Tool Lanes}
% =============================================================================

\begin{frame}{Choose One Lane}
  \begin{tabular}{p{0.12\linewidth}p{0.28\linewidth}p{0.48\linewidth}}
    \toprule
    Lane & Tool & Best fit \\
    \midrule
    A1 & Codex CLI & terminal-first workflows and GitHub tasks \\
    A2 & Codex app & project threads, long goals, visual review \\
    B1 & Claude Code CLI & terminal coding and explicit subagents \\
    B2 & Claude Code app & app-based code sessions \\
    C & Cursor & IDE-first coding, rules, background agents \\
    \bottomrule
  \end{tabular}
  \vspace{0.5em}
  \handoff{\artifact{materials/tool\_tracks.md}}
\end{frame}

\begin{frame}{Codex Lane}
  \begin{itemize}
    \item Use \artifact{AGENTS.md} and project-specific instructions.
    \item CLI and app share the same project files.
    \item Good for issue-driven implementation, long goals, and PR workflows.
    \item Verify with commands and attach evidence to PRs.
  \end{itemize}
  \handoff{\artifact{examples/codex/}, \artifact{examples/codex-app/}}
\end{frame}

\begin{frame}{Claude Code Lane}
  \begin{itemize}
    \item Use \artifact{CLAUDE.md} plus shared \artifact{AGENTS.md} when collaborating.
    \item Skills, subagents, and hooks are useful for research review loops.
    \item Keep project settings in the project, user secrets outside the repo.
    \item Ask the reviewer role for blockers before fixes.
  \end{itemize}
  \handoff{\artifact{examples/claude/}, \artifact{examples/claude-app/}}
\end{frame}

\begin{frame}{Cursor Lane}
  \begin{itemize}
    \item Use rules for persistent project context.
    \item Use IDE review panes and background agents for branch work.
    \item Skills and subagents can mirror the same economics harness pattern.
    \item Verify UI labels and hook support against the installed Cursor version.
  \end{itemize}
  \handoff{\artifact{examples/cursor/}}
\end{frame}

\begin{frame}{Version-Sensitive Tool Claims}
  \begin{itemize}
    \item Tool surfaces change quickly.
    \item Do not memorize UI labels from slides.
    \item Use \artifact{docs/sources.md} for official docs.
    \item If a feature is missing, fall back to files and GitHub issues.
  \end{itemize}
\end{frame}

% =============================================================================
\section{Orchestration}
% =============================================================================

\begin{frame}{GitHub As Control Plane}
  \begin{itemize}
    \item Issue: task card with scope and acceptance criteria.
    \item Branch: isolated implementation.
    \item PR: reviewable diff plus verification evidence.
    \item Merge: only after evidence and human review.
  \end{itemize}
  \handoff{\artifact{orchestration/cloud\_agent\_issue\_template.md}}
\end{frame}

\begin{frame}[fragile]{Issue Template For Agents}
  \begin{lstlisting}[style=markdown]
Title: Add Card-Krueger robustness check

Allowed files:
- examples/card-krueger/src/did_analysis.py
- examples/card-krueger/tests/test_did_analysis.py

Acceptance criteria:
- baseline output remains unchanged;
- robustness result is documented;
- pytest passes;
- synthetic-data caveat stays visible.
  \end{lstlisting}
\end{frame}

\begin{frame}{Cloud Agents And PR Evidence}
  \begin{itemize}
    \item Give cloud agents small, independent issues.
    \item Require branch name, allowed files, and verification output.
    \item Do not accept "works" without commands and output.
    \item Use reviewer agents on the PR, not in the same implementation thread.
  \end{itemize}
\end{frame}

\begin{frame}{Swarm Planning}
  \begin{tabular}{p{0.23\linewidth}p{0.38\linewidth}p{0.25\linewidth}}
    \toprule
    Stream & Task & Depends on \\
    \midrule
    Data map & source, variables, caveats & none \\
    Estimation & baseline and tests & data map \\
    Literature & citations and claim checks & none \\
    Review & verification and blockers & PRs \\
    \bottomrule
  \end{tabular}
  \vspace{0.5em}
  Use labels: \artifact{parallel}, \artifact{sequential}, \artifact{reviewer}.
  \handoff{\artifact{orchestration/card\_krueger\_swarm.md}}
\end{frame}

\begin{frame}{When Cloud Tools Fail}
  \begin{itemize}
    \item Record the failure exactly: permission, missing dependency, usage limit, or build error.
    \item Use local fallback only for the blocked slice.
    \item Preserve small commits and PR evidence.
    \item Return to cloud agents on the next independent issue.
  \end{itemize}
\end{frame}

% =============================================================================
\section{Sprint And Verification}
% =============================================================================

\begin{frame}{Research Sprint}
  \begin{itemize}
    \item Pick one artifact: memo section, test, table, data map, or review log.
    \item Give the agent one file set and one output.
    \item Stop after the first useful diff.
    \item Verify before expanding scope.
  \end{itemize}
  \handoff{\artifact{notes/orchestration\_log\_template.md}}
\end{frame}

\begin{frame}{Verification Stack}
  \begin{itemize}
    \item Unit tests for code paths.
    \item Replication checks for outputs and tables.
    \item Referee-style review for identification and overclaiming.
    \item Playwright or screenshots for browser evidence.
    \item Human sign-off for interpretation.
  \end{itemize}
  \handoff{\artifact{docs/research\_article\_harness.md}}
\end{frame}

\begin{frame}[fragile]{Verification Prompt}
  \begin{lstlisting}[style=markdown]
Act as an economics referee. Review only
examples/card-krueger/docs/research_design_memo.md.

Return blockers first:
- identification;
- data source;
- sample definition;
- outcome units;
- overclaiming.

Do not edit files.
  \end{lstlisting}
\end{frame}

\begin{frame}{Adoption Checklist}
  \begin{itemize}
    \item Choose one real project.
    \item Add an instruction file and privacy boundary.
    \item Add one data source map.
    \item Add one verification command or checklist.
    \item Create one issue with acceptance criteria.
    \item Run one reviewer pass before trusting output.
  \end{itemize}
\end{frame}

\begin{frame}{Close}
  \begin{block}{Takeaway}
    Agentic coding is useful for economists when delegation is bounded by
    context, acceptance criteria, and verification.
  \end{block}
  \begin{itemize}
    \item Start small.
    \item Keep the human in charge of research claims.
    \item Make every agent task leave evidence.
  \end{itemize}
\end{frame}

\end{document}
